Cognitive models cast psychological theories in probabilistic form, specifying how unobservable mental quantities (e.g., processing speed) give rise to observable behavior (e.g., reaction times).
These models are typically \textit{generative}; that is, they define a recipe for simulating realistic behavior from any given configuration of the underlying constructs.
Traditionally, this forward machinery is handcrafted anew for each experiment: one model specification per task, simulated and analyzed on its own.
A recent development breaks with this convention by moving one level up, from models to \textit{meta-models}: single systems trained over an entire space of tasks or model specifications, which then generalize across that space.
The most prominent example to date is Centaur \citep{binz2025centaur}, a language model fine-tuned on trial-by-trial data from hundreds of psychological experiments that predicts human behavior in held-out tasks with remarkable fidelity.
Viewed through this lens, Centaur is a meta-model of the \textit{forward} direction: instead of implementing any particular cognitive model, it subsumes the behavior-generating machinery of many task-specific models into one system.

Yet the scientific payoff of cognitive modeling lies almost entirely in the opposite, \textit{inverse} direction: given observed behavior, what can we conclude about the latent constructs that produced it?

Because this mapping does not typically preserve all information about $\rvtheta$, inference needs to represent uncertainty via the posterior $p(\rvtheta \mid \rvy)$. Since posteriors of cognitive models are rarely tractable, inference tends to rely on iterative Monte Carlo methods, which can take up to multiple days \citep[e.g.,][]{stricklandRacingRemember2018, mileticParameterRecovery2017}.
Here, the meta-modeling revolution has not yet arrived. Centaur itself offers no help: it exposes no interpretable latent constructs, and it cannot tell a researcher what an observed pattern of responses implies about processing speed, response caution, or memory capacity.
The inverse counterpart---a meta-model that maps behavior to posterior distributions over the parameters of interpretable cognitive models, across an entire space of such models---does not yet exist.
This asymmetry is consequential, because measurement, theory testing, and computational phenotyping all operate in the inverse direction.

The natural starting point for building such an inverse meta-model is neural simulation-based inference \citep[SBI;][]{cranmer2020frontier, deistler2025simulation}, which employs neural networks trained on simulated data to avoid iterative estimation. In particular, \textit{amortized methods} \citep{burkner2023some} learn from simulated parameter–observation pairs $(\rvtheta, \rvy)$, enabling rapid sampling from the posterior $p(\rvtheta \mid \rvy^*)$ for any new observation $\rvy^*$ without retraining.
Cognitive models are, in many ways, ideal generators of the required synthetic data: they are easy to implement, fast to simulate (often in the millisecond range), and low-dimensional.
Accordingly, SBI applications in cognitive modeling abound \citep[e.g.,][]{radev2020amortized, wieschen2020jumping, fengler2021likelihood, boelts2022flexible, von2022mental, ghaderi2023general, sokratous2023ask, schumacher2023superstat, schumacher2024validation, rmus2024artificial, hato2025, schaefer2025conflictModels, scholten2026diffusionLens, kvam2024using, kvam2025comparing, belov2026supervised, wu2026testing}.
Standard amortization, however, operates one level below the meta-model: the network is amortized \textit{within} a single, fixed model specification, so its scope ends where that specification does.

This mismatch between per-model amortization and modeling practice is not a corner case; it is the norm.
In practice, cognitive inference rarely involves fitting only a single model.
Instead, we fit multiple candidate models, run many iterative cycles of free vs.~fixed parameter configurations, regress parameters on covariates, and generally reconfigure the model as theory and data negotiate \citep{voss2019sequential, evans2019assessing, stevensonBayesianHierarchical2026}.
This dynamic process can be packaged into principled \textit{Bayesian workflows} \citep{gelman2020bayesian, schad2021toward, li2024amortized}.
For example, when modeling response times with Sequential Sampling Models \citep[SSM;][]{smith2025ssmBible}, a researcher might compare variants with a fixed decision threshold to ones with a collapsing threshold that decreases over time, repeatedly refitting the model while constraining or freeing parameters, or testing whether an experimental manipulation affects a single parameter (e.g., drift rate) or multiple components of the decision process \citep{von2022mental}.
Under per-model amortization, every such reconfiguration demands re-training the network, steadily eroding the very amortization benefits that made SBI attractive in the first place.
What the workflow demands is \textit{meta-amortization}: amortizing not only over datasets of a given model, but over the space of model specifications itself.

The current work delivers such a meta-amortized inference engine, building on a stream in SBI that extends amortization across models \citep{chang2025amortized, gloeckler2024all, elsemuller2023sensitivity, schroder2023simultaneous}:
training a single inference network that remains valid across many, potentially infinitely many, structurally similar models.
We dub our architecture the \textbf{CogFormer}, directly inspired by the \textbf{SimFormer} \citep{gloeckler2024all}.
CogFormer can adapt to any configuration within a model collection's combinatorial parameter space and infer its components accordingly, without substantially compromising the quality of inference when compared to single-model posterior approximation.
In this work, we taxonomize the scope of amortization in cognitive modeling and demonstrate compelling results from the first meta-amortized inference engine for cognitive modeling: where Centaur offers a meta-model of how behavior is generated, CogFormer offers its inverse---a meta-model of how latent cognition is recovered from behavior.

\section{Related work}
\label{sec:related}

Meta-amortization has earlier roots in variational settings. For example, MetaVAE \citep{wu2020meta} proposed doubly amortized inference across a family of related probabilistic models via a MetaELBO objective, but it was not designed for flexible posteriors in SBI. Within SBI, earlier work considered simultaneous inference both over discrete model indices \citep{radev2020amortized} and over discrete models and continuous parameters for compositional simulators \citep{schroder2023simultaneous}.

A separate line of work moves toward multi-query SBI, in which a single amortized artifact can answer many conditional inference queries at test time. For instance, SimFormer \citep{gloeckler2024all} trains a transformer-based diffusion model on the joint distribution of parameters and observations and uses masking and attention structure to sample arbitrary conditionals, including posteriors and likelihoods. OneFlowSBI \citep{nautiyal2026oneflowsbi} pursues a similar goal with flow matching: it learns a vector field over the joint state space and realizes posterior, likelihood, and mixed conditionals. \citet{reuter2025can} amortizes over different statistical model families (GLMs, latent factor models) within a single in-context learning framework. These primarily expand query flexibility within a single learned joint model and serve as inspiration for our masking approach. In contrast, we shift the emphasis from many conditionals of a single joint distribution to a flexible posterior estimator that spans a power set of structural model configurations.

A complementary approach pursues amortization across different tasks, such as predictive and posterior estimation. For instance, ACE \citep{chang2025amortized} uses transformers to flexibly reinterpret latent and observed variables at training and test time, enabling both runtime prior injection and continuous and discrete targets via parametric heads. Closest in spirit is SA-ABI \citep{elsemuller2023sensitivity}, which formalizes sources of sensitivity, i.e., priors and data model assumptions, as context variables and uses weight sharing and deep ensembles to reduce refitting in sensitivity analysis. Similarly, the Distribution Transformer \citep{whittle2025distribution} maps prior distributions directly to (time-varying) posteriors via learnable Gaussian mixtures. In our IID cognitive-modeling settings, we take meta-amortization one step further by substantially expanding generalization to combinatorial design spaces and learning an expressive neural estimator without parametric assumptions. Furthermore, we conduct a systematic and multidimensional analysis of amortization gaps.
